\section{From Axis Mappings to Attribute Builders}
\label{sec:axis-mappings-to-attribute-builders}

In this section, we explain how TEPL compiles the complete constructor from
Section~\ref{sec:attribute-synthesis-algorithm-and-guarantees} into a reusable metadata program.
We first introduce Axes IR, which represents operations using symbolic
axis groups, operand bindings, and value semantics. We then explain how
TEPL constructs the metadata program through symbolic axis mapping
and partitioning.

\subsection{Axes IR}
\label{sec:axes-ir}
To reason about an operation across different tensor ranks, TEPL must
describe which input axes share coordinates, which axes appear in the
output, and how output values are computed. Axes IR expresses these
relationships using symbolic axis groups.

Axis attributes are declared in the \code{attrs} block. Each
\code{axis[]} attribute stores an ordered list of dimension positions.
The \code{axes} block uses these attributes with \code{bind} to assign
a symbolic name to a group of input axes. It then uses \code{yield}
to specify the output axis order. The \code{semantics} block describes
how to compute each output element using these axis groups.
Listing~\ref{lst:dot-declaration} illustrates this structure with
generalized dot.

\begin{lstlisting}[language=TEPL,
  caption={Axes IR declaration for generalized dot.},
  label={lst:dot-declaration}]
op dot(lhs: tensor, rhs: tensor) -> tensor {
  attrs {
    lb, rb: axis[] // left and right batch dimensions
    lc, rc: axis[] // left and right contraction dimensions
  }
  axes {
    let lf = ordered_complement(all_axes(lhs), attrs.lb ++ attrs.lc) // left free
    let rf = ordered_complement(all_axes(rhs), attrs.rb ++ attrs.rc) // right free
    bind B { lhs: attrs.lb, rhs: attrs.rb }
    bind K { lhs: attrs.lc, rhs: attrs.rc }
    bind I { lhs: lf }
    bind J { rhs: rf }
    yield B ++ I ++ J
  }
  semantics {
    yield Reduce(K, Product(lhs, rhs))
  }
}
\end{lstlisting}

For Listing~\ref{lst:dot-declaration}, consider operand axis orders
$(i,k_0,b,k_1)$ and $(k_1,j,k_0,b)$.
With \code{lb=[2]}, \code{rb=[3]}, \code{lc=[1,3]}, and \code{rc=[2,0]},
the bindings give $\vec B=(b)$ and $\vec K=(k_0,k_1)$.
Each binding pairs operand positions in list order; paired axes share
a coordinate and must have equal extents.
Removing these positions with \code{ordered_complement} leaves the
free groups $\vec I=(i)$ and $\vec J=(j)$.
Group lengths may vary across matches, including zero.

The statement \code{yield B ++ I ++ J} sets output order $(b,i,j)$.
The \code{semantics} block multiplies input elements and sums over
\code{K}:
\begin{equation}
 T_{b i j}=\sum_{k_0,k_1}
 \mathrm{lhs}_{i k_0 b k_1}\,\mathrm{rhs}_{k_1 j k_0 b}.
\end{equation}

Appendix~\ref{app:operation-semantics} explains how the \code{semantics}
block expresses operation values and input accesses.
The bindings provide the initial axis mappings used below.

\subsection{Symbolic Axis Mapping}
\label{sec:symbolic-axis-mapping}
The symbolic groups in Section~\ref{sec:axes-ir} describe axes within
one operation. When operations are composed, TEPL must track these
axes across tensors and group definitions. A later operation may select
only part of an earlier group, so positions within that group must be
translated into positions in the tensor.

Each \code{bind} statement supplies a position-list expression $p_G$
for a group $\vec G$ in an operand $t$. This mapping satisfies
\begin{equation}
 \gath(\operatorname{Ord}_t,p_G)=\vec G,
 \label{eq:mapping}
\end{equation}
where $\gath$ selects entries from a list at the given positions.
The expression $p_G$ is known from the operation declaration, although
its concrete values may depend on captured attributes.
The \code{yield} statement determines how these groups appear in the
output and therefore provides their output positions.

For example, consider the dot declaration in
Listing~\ref{lst:dot-declaration}, with operand axis orders
$[i_0,i_1,k]$ and $[k,j_0,j_1]$. Suppose the dot contracts $k$ and
has no batch axes. The bindings \code{bind I { lhs: lf }} and
\code{bind J { rhs: rf }} identify the two free-axis groups.
Here \code{lf=[0,1]} and \code{rf=[1,2]}, so the bindings identify
$\vec I=[i_0,i_1]$ and $\vec J=[j_0,j_1]$.
Since \code{yield B ++ I ++ J} places $\vec I$ before $\vec J$,
the result $T$ has
$\operatorname{Ord}_T=[i_0,i_1,j_0,j_1]$.
The groups therefore have positions $p_I=[0,1]$ and $p_J=[2,3]$
in $T$. These are positions in the result, rather than the operand
positions supplied by the bindings.

Now suppose a later operation selects $j_1$ from $\vec J$.
Its position within $\vec J$ is $h=[1]$, but its position in $T$ is
$\gath(p_J,h)=\gath([2,3],[1])=[3]$.
In general, a selection $h$ within a mapped group $\vec G$ has
tensor positions $\gath(p_G,h)$ because
\begin{equation}
 \gath(\gath(\operatorname{Ord}_t,p_G),h)
 =\gath(\operatorname{Ord}_t,\gath(p_G,h)).
 \label{eq:gather-composition}
\end{equation}
Both position lists must be valid for the lists they select from.

TEPL composes these mappings as it tracks axes through the expression
graph. Each selection produces a position-list expression in terms of
the existing mappings. The compiler retains these expressions
symbolically, allowing group lengths and positions to vary across
matches. Evaluating them on match metadata gives the concrete positions
needed for RHS axis attributes.

\subsection{Symbolic Axis Partitioning}
\label{sec:symbolic-axis-partitioning}
Consider the dot reassociation rule
\[
 \operatorname{dot}(\operatorname{dot}(X,Y),Z)
 \;\longrightarrow\;
 \operatorname{dot}(X,\operatorname{dot}(Y,Z)).
\]
On the LHS, let $T=\operatorname{dot}(X,Y)$ and
$O=\operatorname{dot}(T,Z)$. On the RHS, let
$U=\operatorname{dot}(Y,Z)$, so the replacement computes
$\operatorname{dot}(X,U)$.
Changing the graph also changes which axes each dot must contract.
We must determine which contractions can move into $U$ and which
must remain in the final dot, while preserving the original axis pairs.

For clarity, consider dots without batch axes.
The first LHS dot contracts group $\vec K$ between $X$ and $Y$.
Its output $T$ contains the free groups $\vec I$ from $X$ and
$\vec J$ from $Y$, in that order.
The second dot selects contraction axes from $T$ using the position
list $c_L$ and pairs them with positions $c_R$ in $Z$.
These lists pair axes entry by entry: slot $s$ pairs the axis at
$c_L[s]$ in $T$ with the axis at $c_R[s]$ in $Z$.
The selected axes may come from both $\vec I$ and $\vec J$.
TEPL must separate them by producer group and retain their partners
in $Z$.

For example, let
$\operatorname{Ord}_T=[i_0,i_1,j_0,j_1]$, with
$p_I=[0,1]$ and $p_J=[2,3]$.
Suppose $c_L=[3,0]$ and $c_R=[2,0]$.
The second dot selects $[j_1,i_0]$ from $T$:
$j_1$ is paired with position $2$ in $Z$, and $i_0$ with position $0$.
Thus $\vec I$ splits into selected axes $[i_0]$ and remaining axes
$[i_1]$, while $\vec J$ splits into $[j_1]$ and $[j_0]$.
Figure~\ref{fig:axis-partitioning} shows these selections and pairings.

\paperfigure{1}{figures/Symbolic Axis Partitioning (Dot Reassociation).pdf}{Symbolic axis partitioning for unbatched dot reassociation. Splitting selected axes by producer origin preserves their contraction partners. Here $L$ is the selected contraction group, and $F=[f]$ is $Z$'s free group.}{fig:axis-partitioning}{Three panels show the intermediate tensor T with axes i0, i1, j0, j1; selection of positions 3 and 0; partition into selected and remaining I and J subgroups; and preservation of partners at positions 2 and 0 of Z. A table distinguishes positions within producer groups from slots in the contraction list and corresponding partner positions.}

Two kinds of positions are needed for each selected subgroup.
Positions within the producer group identify which axes to select.
Slots in the contraction lists identify their partners.
For instance, $j_1$ is at position $1$ within $\vec J$, but it occupies
slot $0$ in the selection $[j_1,i_0]$.
Position $1$ selects $j_1$ from $\vec J$; slot $0$ selects its partner
position $c_R[0]=2$ in $Z$.
Neither position can replace the other.

For each producer group $\vec G$, let $v_G$ list the selected positions
within $\vec G$, and let $u_G$ list their slots in $c_L$.
The mapping $p_G$ from Section~\ref{sec:symbolic-axis-mapping}
provides the positions of all axes in $\vec G$ within $T$.
The compiler compares $c_L$ with $p_G$:
\begin{equation}
 (u_G,v_G)=\mpo(c_L,p_G),\qquad
 z_G=\gath(c_R,u_G).
 \label{eq:symbolic-partition-selectors}
\end{equation}
Here $\mpo$ finds entries common to the two lists and returns their
positions in each list, following the order of $c_L$.
A common entry identifies one tensor axis:
$c_L[u_G[s]]=p_G[v_G[s]]$.
The list $z_G$ then gives its partner positions in $Z$.
Appendix~\ref{app:candidate-validation-and-soundness} gives the formal
definition.

For the running example, these lists are:
\begin{center}
\small
\begin{tabular}{@{}lllll@{}}\toprule
Group & Selected axes & $v_G$: in group & $u_G$: pairing slots
      & $z_G$: in $Z$\\\midrule
$\vec I$ & $[i_0]$ & $[0]$ & $[1]$ & $[0]$\\
$\vec J$ & $[j_1]$ & $[1]$ & $[0]$ & $[2]$\\\bottomrule
\end{tabular}
\end{center}
The selected subgroup is $\gath(\vec G,v_G)$.
Its positions in $T$ are $\gath(p_G,v_G)$, and its partners occupy
$z_G$ in $Z$. The ordered complement of $v_G$ within the group
selects the remaining axes. Selected axes follow contraction-list
order; remaining axes preserve their original group order.

These mappings determine how to distribute the contractions on the RHS.
Since $\vec J$ originates in $Y$, the selected subgroup
$\vec J_{\mathrm{sel}}$ can contract with its partners in $Z$ inside
$U=\operatorname{dot}(Y,Z)$.
TEPL uses $v_J$ to select these axes through the mapping of $\vec J$
into $Y$, and $z_J$ to select their partners in $Z$.
In the example, this contracts $j_1$ with position $2$ in $Z$.

The selected axes from $\vec I$ originate in $X$, which is absent
from $U$. Their partners in $Z$ must therefore survive in $U$.
Group $\vec K$ from $Y$ must also survive so that it can later
contract with $X$.
The final dot contracts $\vec K$ and pairs
$\vec I_{\mathrm{sel}}$ with its retained partners.
The mappings through $U$ provide their new positions.
The remaining groups $\vec I_{\mathrm{rem}}$,
$\vec J_{\mathrm{rem}}$, and the free axes of $Z$ form the output.
Partitioning therefore supplies both the axis subsets and the pairing
information needed to construct the RHS attributes.

\subsection{Symbolic Attribute Builder Construction}
\label{sec:symbolic-attribute-builder-construction}
The mappings and partitions from the preceding sections identify the
axes needed by each RHS operation. The compiler turns these results
into expressions that compute its axis attributes. Together with the
feasibility checks, these expressions form a \emph{metadata program}:
a program that computes on captured attributes, ranks, and shapes.

Continuing the reassociation example, $U=\operatorname{dot}(Y,Z)$
contracts $\vec J_{\mathrm{sel}}$ with its partners in $Z$.
Let $p_J^Y$ map $\vec J$ to its positions in $Y$.
The selector $v_J$ identifies the selected axes within $\vec J$,
while $u_J$ identifies their slots in the original contraction lists.
The two contraction attributes for $U$ are therefore
\begin{equation}
 c_L^U=\gath(p_J^Y,v_J),\qquad
 c_R^U=\gath(c_R,u_J).
 \label{eq:builder-contraction-example}
\end{equation}
For the operand order $\operatorname{Ord}_Y=[k,j_0,j_1]$ from
Section~\ref{sec:symbolic-axis-mapping}, $p_J^Y=[1,2]$.
With $v_J=[1]$, $u_J=[0]$, and $c_R=[2,0]$, the expressions give
$c_L^U=[2]$ and $c_R^U=[2]$. These lists tell the RHS dot to contract
$j_1$ in $Y$ with its original partner in $Z$.

At compilation time, TEPL records these expressions and their
dependencies on source metadata. For example, $v_J$ and $u_J$ are
computed by the $\mpo$ expression from
Equation~\eqref{eq:symbolic-partition-selectors}.
Shared computations are represented once and reused across attributes.
The compiler constructs this expression graph once per rule, without
enumerating concrete axis lists. At each match, evaluating the graph
supplies the current position lists and computes the RHS attributes.
The same expressions therefore handle different ranks and group lengths.

The compiler applies this approach to all three passes in
Section~\ref{sec:attribute-synthesis-algorithm-and-guarantees}:
determining retained axes, checking order, and constructing attributes.
Write $E_{\mathrm{check}}$ for the generated feasibility check and
$E_{\mathrm{build}}$ for the generated attribute builder.
The notation $\lbrack\!\lbrack E\rbrack\!\rbrack_m$ denotes evaluating
program $E$ on match metadata $m$.
The check runs first; the builder is evaluated only when it succeeds.

Several attribute orders may describe the same valid axis pairing.
To state exact agreement with the mathematical builder, both use the
same ordering convention $\tau$: unconstrained batch and contraction
axes follow the left operand's order.
Let $\operatorname{build}_{\tau}$ denote the builder from
Section~\ref{sec:constructing-axis-attributes} with this convention.

\begin{theorem}[Symbolic compilation correspondence]
\label{thm:symbolic-compilation}
For every rewrite rule satisfying the synthesis restrictions in
Section~\ref{sec:formal-scope-and-guarantees}, compilation constructs
a finite metadata program $(E_{\mathrm{check}},E_{\mathrm{build}})$
independently of concrete ranks and the written guard.
For every legal metadata instance $m$ in that section's completeness domain,
\begin{align}
 \lbrack\!\lbrack E_{\mathrm{check}}\rbrack\!\rbrack_m&=\operatorname{check}(m),
 \label{eq:emitted-check}\\
 \operatorname{check}(m)&\Rightarrow
 \lbrack\!\lbrack E_{\mathrm{build}}\rbrack\!\rbrack_m=\operatorname{build}_{\tau}(m).
 \label{eq:emitted-build}
\end{align}
\end{theorem}
The compiled check makes the same feasibility decision as the
mathematical algorithm. When it succeeds, the compiled builder produces
the same attributes under $\tau$.
Together with Theorem~\ref{thm:complete-synthesis}, this gives a complete
compiled synthesis procedure within the stated domain.

The proof follows the three passes and shows that the generated
expressions preserve their axis sets, ordering constraints, and position
queries. Appendix~\ref{app:symbolic-program-construction-and-correspondence}
gives the proof, including termination of program construction.
Appendix~\ref{app:batched-reassociation} gives a complete reassociation
builder under explicit sufficient conditions.

\subsection{Compilation and Validation}
\label{sec:compilation-and-validation}
The compiled programs implement Section~\ref{sec:synthesis-and-execution-contract}'s execution
contract. Before accepting them, the compiler submits the feasibility
check and completed attribute builder to conditional validation
(Section~\ref{sec:candidate-validation}). Validation establishes that the
check is defined using source metadata, then proves the builder's
obligations under the applicability condition. Other construction
obligations remain goals to prove.

For sound schemas, validation also checks the schema conditions and
agreement of supplied fields with constructed fields. Table~\ref{tab:guarantee-map}
distinguishes this soundness result from complete synthesis and compilation
correspondence.
